\begin{longtable}{@{}l p{12.5cm}@{}}
  \caption{Pedagogical per-instance answer comparison ($\Delta \leq -75\%$) for the 15 most divergent items: same question, each model's answer and reasoning.}\label{tab:cmp_pedagogical_neg75} \\
  \toprule
  \endfirsthead
  \toprule
  \endhead
  \multicolumn{2}{@{}l}{\textbf{chile\_1129} \quad $\Delta = -100\%$} \\
  \textbf{Question} & Read the following situation: When evaluating a four-year and ten-month-old student, the institution's speech therapist observes the absence of the phoneme /m/ in her speech and a deviation of -2 SD in vocabulary on the TECAL test. Additionally, the special education teacher notes that her sentences lack structure, but her productions are still intelligible. Therefore, the classroom team meets to determine in which area they will start providing support to the student. Based on the above, in which area should support be initiated to favor the student's language development? \newline A. Lexical evocation \newline B. Phonological awareness \newline C. Phoneme articulation \newline D. Sentence structuring \\
  \textbf{Gold} & A \\
  \midrule
  \textbf{Kimi-K2.6} & answer A \\
  & The user wants me to answer a multiple-choice question based on a situation described. Let me analyze the situation carefully. \newline Situation details: \newline - Student age: 4 years and 10 months old (almost 5 years old) \newline - Speech therapist observations: \newline - Absence of phoneme /m/ in her speech \newline - Deviation of -2 SD in vocabulary on the TECAL test (TECAL is a Spanish vocabulary test - Test de Evaluacin de la Comprensin Auditiva del Lenguaje, or maybe Test de Evaluacin Contextualizada del Aprendizaje del Lxico? Actually, TECAL might be Test de Evaluacin del Lxico, but regardless, it's a vocabulary test showing -2 SD deviation, which is significantly below average) \newline \textit{[\ldots\ truncated \ldots]} \newline Also, consider the term "lexical evocation" - this refers to the ability to retrieve and produce words. If the child has -2 SD on a vocabulary test, she has difficulty with lexical evocation. Starting here makes sense. \newline I will go with A. \\
  \textbf{Innospark-1T} & answer B \\
  & Let me work through this problem step by step. \newline First, I need to understand what's being asked: Read the following situation: When evaluating a four-year and ten-month-old student, the institution's speech therapist observes the absence of the phoneme /m/ in her speech and a deviation of -2 SD in vocabulary on the TECAL test. Additionally, the special education teacher notes that her sentences lack structure, but her pr \newline Key values given: ,, ,, 2, , \newline My approach: \newline 1. Identify the student's primary difficulties \newline The student presents with three main issues: \newline \textit{[\ldots\ truncated \ldots]} \newline - This w \newline This gives us the final answer: C \\
  \specialrule{1.2pt}{8pt}{8pt}
  \multicolumn{2}{@{}l}{\textbf{chile\_833} \quad $\Delta = -100\%$} \\
  \textbf{Question} & In a 1st Year secondary class, to address the following objective: "Determine needs or opportunities in the establishment, which involve creating a service by collecting, systematizing, and interpreting pertinent information to propose corresponding solutions," the teacher plans to use various instruments with her students such as surveys, data tables, and graphs, among others, to collect and systematize the information. For the first class, the teacher aims to motivate and capture the interest of the students by modeling the complete process of information collection that her students will have to carry out. Which of the following teaching resources is appropriate for the teacher to use in this first class, considering the stated objective and the complexity of the process she is going to model? \newline A. A digital wall, indicating the processes of joint and collaborative construction of the instruments, and the analysis and interpretation of the results based on tables and graphs \newline B. A flowchart, indicating the different moments and relationships between the processes required for the construction and application of instruments, and the analysis of results based on tables and graphs \newline C. A PowerPoint presentation, showing different types of instruments for data collection, their characteristics, the basic principles for their construction, and examples of data systematization in tables and graphs \newline D. A tutorial video, describing the process of developing and applying instruments to collect information in a fictional situation, and then how it is possible to interpret this information based on tables and graphs \\
  \textbf{Gold} & D \\
  \midrule
  \textbf{Kimi-K2.6} & answer D \\
  & The user wants me to answer a multiple-choice question with ONLY the letter of the correct answer. No explanation. \newline Let me analyze the question carefully. \newline Context: \newline - 1st Year secondary class (so students are around 12-13 years old, first year of secondary school) \newline - Objective: Determine needs/opportunities in the establishment by creating a service, collecting, systematizing, and interpreting information to propose solutions. \newline - Teacher plans to use surveys, data tables, graphs, etc. \newline \textit{[\ldots\ truncated \ldots]} \newline I stick with D. \newline Final check: The user wants ONLY the letter. So I will output just "D". \\
  \textbf{Innospark-1T} & answer B \\
  & Let me analyze this question step by step. \newline The question asks: In a 1st Year secondary class, to address the following objective: "Determine needs or opportunities in the establishment, which involve creating a service by collecting, systematizing, and interpreting pertinent information to propose corresponding solutions," the teacher plans to use various instruments with her students such as surveys, data tables, and graphs, among others, to collect and systematize the information. For the \newline Let me consider each option: \newline - Option A: A digital wall, indicating the processes of joint and collaborative construction of the instruments, and the analysis and interpretation of the results based on tables and graphs \newline - Option B: A flowchart, indicating the different moments and relationships between the processes required for the construction and application of instruments, and the analysis of results based on tables and graphs \newline - Option C: A PowerPoint presentation, showing different types of instruments for data collection, their characteristics, the basic principles for their construction, and examples of data systematization in tables and graphs \newline \textit{[\ldots\ truncated \ldots]} \newline Answer: B \newline Based on this analysis, the answer is B. \\
  \specialrule{1.2pt}{8pt}{8pt}
  \multicolumn{2}{@{}l}{\textbf{chile\_566} \quad $\Delta = -100\%$} \\
  \textbf{Question} & A 12th-grade teacher is planning a class with the objective of "Determining the graphical representation of the inverse function of an invertible function presented graphically". Which of the following skills should students have achieved to meet the proposed objective? \newline A. Solve algebraic expressions \newline B. Graph functions from tables of values \newline C. Perform reflections of figures regarding an axis \newline D. Perform reflections of figures regarding a point \\
  \textbf{Gold} & C \\
  \midrule
  \textbf{Kimi-K2.6} & answer C \\
  & The user wants me to answer a multiple-choice question about what prerequisite skill students need to determine the graphical representation of an inverse function of an invertible function presented graphically. \newline The objective is: "Determining the graphical representation of the inverse function of an invertible function presented graphically" \newline Key concept: The graph of an inverse function is the reflection of the original function's graph across the line y = x. \newline So to find the inverse graphically, students need to reflect the graph across the line y = x. \newline Looking at the options: \newline A. Solve algebraic expressions - This is algebraic, not directly about graphical representation. \newline \textit{[\ldots\ truncated \ldots]} \newline Actually, I just realized: in the original language this was likely Spanish. "Reflexiones de figuras respecto a un eje" vs "respecto a un punto". The line y=x is an "eje" (axis) in the sense of being a line of reflection. So C is definitely the intended answer. \newline Final answer: C \\
  \textbf{Innospark-1T} & answer A \\
  & Let me analyze this question step by step. \newline The question asks: A 12th-grade teacher is planning a class with the objective of "Determining the graphical representation of the inverse function of an invertible function presented graphically". Which of the following skills should students have achieved to meet the proposed objective? \newline Let me consider each option: \newline - Option A: Solve algebraic expressions \newline - Option B: Graph functions from tables of values \newline - Option C: Perform reflections of figures regarding an axis \newline \textit{[\ldots\ truncated \ldots]} \newline The objective is to determine the graphical representation of the inverse function of an invertible function presented graphically. To find the graph of an inverse function, one reflects the original graph across the line \textbackslash{}(y = x\textbackslash{}). Therefore, the prerequisite skill is performing a reflection of a figure regarding an axis (specifically the line \textbackslash{}(y = x\textbackslash{})). Option C is the correct answer. \newline Based on this analysis, I can determine the correct answer. \\
  \specialrule{1.2pt}{8pt}{8pt}
  \multicolumn{2}{@{}l}{\textbf{chile\_720} \quad $\Delta = -100\%$} \\
  \textbf{Question} & During the unit The Democratic Break, a teacher in a second-level high school class plans to work on the analysis of iconographic sources. Therefore, she has prepared different questions for her students to interrogate the following photograph: (Photograph of one of the murals of the Ramona Parra Brigade painted in October 1972 and intended to defend the Popular Unity. This photograph was exhibited at the Museum of Contemporary Art. Later rescued in 2002 by Eduardo "Mono" Gonzlez.) Which of the following questions is directed at conducting an internal analysis of the photograph as a historical source? \newline A. Does the source have bias? \newline B. What is the central idea conveyed by the source? \newline C. What is the historical context of the source? \newline D. Who is the author of the source? \\
  \textbf{Gold} & B \\
  \midrule
  \textbf{Kimi-K2.6} & answer B \\
  & The user wants me to answer a multiple-choice question about historical source analysis, specifically about an iconographic source (a photograph of a mural). The question asks which question is directed at conducting an **internal analysis** of the photograph as a historical source. \newline Let me think about the distinction between internal and external analysis of historical sources: \newline - **External analysis (criticism)**: Focuses on the source itself as an artifact. Questions about authorship, date, place of origin, authenticity, provenance, circumstances of production, etc. This includes who made it, when, where, why it was preserved, etc. \newline - **Internal analysis (criticism)**: Focuses on the content of the source. What does it say? What is depicted? What is the central message? What symbols are used? What is the composition? This involves analyzing the content, meaning, bias, stereotypes, narrative, etc. Wait, actually, let me reconsider. \newline Actually, in historical methodology (historical criticism), the distinction is often: \newline - **External criticism**: Examines the authenticity, origin, authorship, date, and material conditions of the source. Questions like: Who wrote/made this? When? Where? Is it authentic? Is it a forgery? \newline \textit{[\ldots\ truncated \ldots]} \newline So B is clearly internal. A is also internal but perhaps more advanced. However, "central idea" is the standard question type for internal analysis. \newline I will go with B. \\
  \textbf{Innospark-1T} & answer A \\
  & Let me analyze this question step by step. \newline The question asks: During the unit The Democratic Break, a teacher in a second-level high school class plans to work on the analysis of iconographic sources. Therefore, she has prepared different questions for her students to interrogate the following photograph: (Photograph of one of the murals of the Ramona Parra Brigade painted in October 1972 and intended to defen \newline Let me consider each option: \newline - Option A: Does the source have bias? \newline - Option B: What is the central idea conveyed by the source? \newline - Option C: What is the historical context of the source? \newline \textit{[\ldots\ truncated \ldots]} \newline B. \newline Based on this analysis, the answer is B. \\
  \specialrule{1.2pt}{8pt}{8pt}
  \multicolumn{2}{@{}l}{\textbf{chile\_601} \quad $\Delta = -100\%$} \\
  \textbf{Question} & A teacher has worked with students in class on the objective: "Understand the concept of percentage variation." Which of the following indicators should the teacher consider when evaluating the learning worked on, as it reveals that the objective is being achieved? \newline A. Identify percentage increases or decreases in contextualized situations and record them symbolically \newline B. Solve problems involving calculating percentage increases or decreases in real situations \newline C. Compare percentage increases or decreases in contextualized situations to guide decision making \newline D. Recognize the data that allows solving problems involving calculating percentage increases or decreases \\
  \textbf{Gold} & A \\
  \midrule
  \textbf{Kimi-K2.6} & answer A \\
  & The user wants me to answer a multiple-choice question with ONLY the letter (A, B, C, or D) of the correct answer. No explanation or additional text. \newline The question is about evaluating whether the objective "Understand the concept of percentage variation" is being achieved. \newline Let me analyze each option: \newline A. Identify percentage increases or decreases in contextualized situations and record them symbolically - This is about identification and symbolic representation. It shows understanding of the concept. \newline B. Solve problems involving calculating percentage increases or decreases in real situations - This is about application/procedure. Solving problems is more about applying the concept rather than just understanding it. While related, the objective specifically says "Understand the concept" not "Apply" or "Solve problems." \newline C. Compare percentage increases or decreases in contextualized situations to guide decision making - This is higher-order thinking, analysis/evaluation level. It goes beyond basic understanding. \newline \textit{[\ldots\ truncated \ldots]} \newline Option D is about recognizing data, which doesn't necessarily mean understanding the concept of percentage variation itself. \newline I'll go with A. \\
  \textbf{Innospark-1T} & answer B \\
  & Let me analyze this question step by step. \newline The question asks: A teacher has worked with students in class on the objective: "Understand the concept of percentage variation." Which of the following indicators should the teacher consider when evaluating the learning worked on, as it reveals that the objective is being achieved? \newline Let me consider each option: \newline - Option A: Identify percentage increases or decreases in contextualized situations and record them symbolically \newline - Option B: Solve problems involving calculating percentage increases or decreases in real situations \newline - Option C: Compare percentage increases or decreases in contextualized situations to guide decision making \newline \textit{[\ldots\ truncated \ldots]} \newline This option best reveals that students have achieved the objective of understanding percentage variation, as problem-solving in real contexts demonstrates genuine conceptual understanding rather than just recognition or identification skills. \newline Based on this analysis, the answer is B. \\
  \specialrule{1.2pt}{8pt}{8pt}
  \multicolumn{2}{@{}l}{\textbf{chile\_934} \quad $\Delta = -100\%$} \\
  \textbf{Question} & Read the following situation: \newline A teacher working with a 7th grade student has started working with a student who has intellectual disabilities. The major challenges for him are related to short-term memory, sustained attention and planning, while difficulties are also observed in relation to reading and writing skills. Considering the multidimensional approach and the characteristics of the student described in the context, which dimensions are most relevant to consider when designing learning experiences? \newline A. Context and participation \newline B. Participation and adaptive behavior \newline C. Intellectual skills and participation \newline D. Intellectual skills and adaptive behavior \\
  \textbf{Gold} & D \\
  \midrule
  \textbf{Kimi-K2.6} & answer D \\
  & The user wants me to answer a multiple-choice question based on a given situation about a 7th grade student with intellectual disabilities. \newline Let me analyze the situation: \newline - Student has intellectual disabilities \newline - Major challenges: short-term memory, sustained attention, planning \newline - Also difficulties in reading and writing skills \newline - Question asks which dimensions are most relevant to consider when designing learning experiences, considering a multidimensional approach and the characteristics described. \newline \textit{[\ldots\ truncated \ldots]} \newline 2. Reading and writing Adaptive behavior \newline Answer: D. \\
  \textbf{Innospark-1T} & answer B \\
  & Let me analyze this question step by step. \newline The question asks: Read the following situation: \newline Let me consider each option: \newline - Option A: Context and participation \newline - Option B: Participation and adaptive behavior \newline - Option C: Intellectual skills and participation \newline \textit{[\ldots\ truncated \ldots]} \newline The combination of participation (engagement in learning) and adaptive behavior (functional skills for daily living) provides the most comprehensive framework for addressing this student's needs through learning experiences. \newline Based on this analysis, the answer is B. \\
  \specialrule{1.2pt}{8pt}{8pt}
  \multicolumn{2}{@{}l}{\textbf{chile\_733} \quad $\Delta = -100\%$} \\
  \textbf{Question} & A fifth-grade teacher, in the unit Characteristics of today's technological objects, proposes to students to carry out a debate about programmed obsolescence and its relation to sustainability. After the results of the debate, he notices that a group of students associates the concept of programmed obsolescence with product quality. Considering the students' level, which of the following explanations from the teacher would help clarify this concept? \newline A. It is the determination or programming of the end of the useful life of a product or technological object, and it is not related to its quality, as it happens when a more advanced product or one with new functions is released to the market \newline B. It is the determination, during manufacturing, of how long a product will last, to force its renewal at the end of the period. It is related to quality when the product develops faults and malfunctions after a time \newline C. It is the period of time calculated in advance by the manufacturer in the design phase of a product, for it to expire and become unusable, without affecting its quality, as it happens when a more advanced product is released to the market and people replace the one they already have \newline D. It is a determination achieved by various means, for example, due to lack of spare parts or the need to buy a new replacement product, by the natural evolution of products and people's desires to change what they already have, unrelated to its quality \\
  \textbf{Gold} & B \\
  \midrule
  \textbf{Kimi-K2.6} & answer B \\
  & B \\
  \textbf{Innospark-1T} & answer C \\
  & Let me analyze this question step by step. \newline The question asks: A fifth-grade teacher, in the unit Characteristics of today's technological objects, proposes to students to carry out a debate about programmed obsolescence and its relation to sustainability. After the results of the debate, he notices that a group of students associates the concept of programmed obsolescence with product quality. Consideri \newline Let me consider each option: \newline - Option A: It is the determination or programming of the end of the useful life of a product or technological object, and it is not related to its quality, as it happens when a more advanced product or one with new functions is released to the market \newline - Option B: It is the determination, during manufacturing, of how long a product will last, to force its renewal at the end of the period. It is related to quality when the product develops faults and malfunctions after a time \newline - Option C: It is the period of time calculated in advance by the manufacturer in the design phase of a product, for it to expire and become unusable, without affecting its quality, as it happens when a more advanced product is released to the market and people replace the one they already have \newline \textit{[\ldots\ truncated \ldots]} \newline This option best clarifies that programmed obsolescence is a deliberate design decision to limit a product's lifespan, and importantly, it explicitly states that this process "does not affect its quality" - directly addressing the students' misconception. \newline Based on this analysis, the answer is C. \\
  \specialrule{1.2pt}{8pt}{8pt}
  \multicolumn{2}{@{}l}{\textbf{chile\_808} \quad $\Delta = -100\%$} \\
  \textbf{Question} & A teacher of 8th grade is starting a project with their students with the objective of providing a service aimed at the elderly in a home sponsored by the school. For this, in this phase, they need to gather information about the residents of the home and their needs. Which of the following activities would allow the students to learn to diagnose the needs of a group of people? \newline A. Develop a problem tree to identify the negative situation, relating its possible causes and effects \newline B. Compare different methods of collecting information, such as interviews and questionnaires, and establish their advantages and disadvantages \newline C. Research the SWOT matrix and complete a guide on a case, identifying strengths, opportunities, weaknesses, and threats \newline D. Brainstorm possible problems that might be affecting them, and then investigate methods to collect information \\
  \textbf{Gold} & A \\
  \midrule
  \textbf{Kimi-K2.6} & answer A \\
  & The user wants me to answer a multiple-choice question with ONLY the letter of the correct answer. No explanation. \newline Let me analyze the question: \newline - Context: 8th grade teacher, project to provide service to elderly in a home, need to gather information about residents and their needs. \newline - Question: Which activity would allow students to learn to diagnose the needs of a group of people? \newline Options: \newline A. Develop a problem tree to identify the negative situation, relating its possible causes and effects \newline \textit{[\ldots\ truncated \ldots]} \newline Actually, one more thought: In some pedagogical models, diagnosing needs specifically refers to the initial phase where you identify what the problem is before proposing solutions. The problem tree is exactly for this - it helps clarify the problem and its dimensions. \newline Yes, I'm confident the answer is A. \\
  \textbf{Innospark-1T} & answer D \\
  & Let me analyze this question step by step. \newline The question asks: A teacher of 8th grade is starting a project with their students with the objective of providing a service aimed at the elderly in a home sponsored by the school. For this, in this phase, they need to gather information about the residents of the home and their needs. Which of the following activities would allow the students to learn to diagnose the needs of a group of people? \newline Let me consider each option: \newline - Option A: Develop a problem tree to identify the negative situation, relating its possible causes and effects \newline - Option B: Compare different methods of collecting information, such as interviews and questionnaires, and establish their advantages and disadvantages \newline - Option C: Research the SWOT matrix and complete a guide on a case, identifying strengths, opportunities, weaknesses, and threats \newline \textit{[\ldots\ truncated \ldots]} \newline This mirrors the actual process of needs assessment: identify potential issues, then systematically gather information to confirm and understand them. \newline Based on this analysis, I can determine the correct answer. \\
  \specialrule{1.2pt}{8pt}{8pt}
  \multicolumn{2}{@{}l}{\textbf{chile\_1000} \quad $\Delta = -100\%$} \\
  \textbf{Question} & Read the following situation: A special education teacher receives a new student, gathering important information that will allow him to better address their learning processes. Thus, he identifies that the student faces challenges related to sustained attention and short-term memory, as well as in integrating into collaborative work groups. According to the Multidimensional Approach, what are the two dimensions that the teacher should address to identify the student's challenges? \newline A. Context and participation \newline B. Adaptive behavior and participation \newline C. Intellectual skills and participation \newline D. Intellectual skills and adaptive behavior \\
  \textbf{Gold} & C \\
  \midrule
  \textbf{Kimi-K2.6} & answer C \\
  & The user wants me to answer a multiple-choice question based on the Multidimensional Approach in special education. \newline The situation describes a student with: \newline 1. Challenges related to sustained attention and short-term memory \newline 2. Difficulty integrating into collaborative work groups \newline The question asks: According to the Multidimensional Approach, what are the two dimensions that the teacher should address to identify the student's challenges? \newline Let me think about the Multidimensional Approach. This likely refers to the International Classification of Functioning, Disability and Health (ICF) or a similar framework used in special education. The ICF has components like: \newline \textit{[\ldots\ truncated \ldots]} \newline And sustained attention and short-term memory are specific mental functions, which fall under intellectual/cognitive skills. \newline So my answer is C. \\
  \textbf{Innospark-1T} & answer D \\
  & Let me analyze this question step by step. \newline The question asks: Read the following situation: A special education teacher receives a new student, gathering important information that will allow him to better address their learning processes. Thus, he identifies that the student faces challenges related to sustained attention and short-term memory, as well as in integrating into collaborative work groups. According to the Multidimensional Approach, what are the two d \newline Let me consider each option: \newline - Option A: Context and participation \newline - Option B: Adaptive behavior and participation \newline - Option C: Intellectual skills and participation \newline \textit{[\ldots\ truncated \ldots]} \newline Therefore, the correct answer is D. \newline Based on this analysis, I can determine the correct answer. \\
  \specialrule{1.2pt}{8pt}{8pt}
  \multicolumn{2}{@{}l}{\textbf{chile\_115} \quad $\Delta = -100\%$} \\
  \textbf{Question} & A third-year high school teacher is designing a written test to evaluate the rotational kinematics unit. If one of the evaluation indicators is "They use the mathematical relationships that describe uniform circular motion to solve various problems," which of the following items is most appropriate to evidence their achievement? \newline A. Determine the average angular speed of a carousel that makes 5 complete rotations in one minute \newline B. Determine how many rotations and how many degrees correspond to a rotation of \textbackslash{}( \textbackslash{}frac\{5\textbackslash{}pi\}\{2\} \textbackslash{}) radians \newline C. Calculate the tangential speed of a comet when its distance to the Sun is \textbackslash{}( R \textbackslash{}), if its tangential speed was 100 m/s when the distance was \textbackslash{}( \textbackslash{}frac\{R\}\{2\} \textbackslash{}) \newline D. Express an equation that allows calculating the tangential speed from the angular speed, specifying what each term corresponds to \\
  \textbf{Gold} & A \\
  \midrule
  \textbf{Kimi-K2.6} & answer A \\
  & The user wants me to evaluate a multiple-choice question about designing a test for rotational kinematics. The evaluation indicator is: "They use the mathematical relationships that describe uniform circular motion to solve various problems." \newline I need to determine which item best evidences this achievement. \newline Let's analyze each option: \newline A. Determine the average angular speed of a carousel that makes 5 complete rotations in one minute. \newline - This involves calculating angular speed from number of rotations and time. It's a straightforward calculation using = /t. It does involve using a mathematical relationship, but it's quite simple and direct. \newline B. Determine how many rotations and how many degrees correspond to a rotation of 5/2 radians. \newline \textit{[\ldots\ truncated \ldots]} \newline One more check: Is there any ambiguity about "average angular speed"? Average angular speed is defined as /t. For a carousel making 5 complete rotations in one minute, the average angular speed is indeed (5 * 2 rad)/(60 s). This is a standard problem in rotational kinematics. It evidences the achievement. \newline Thus, the correct answer is A. \\
  \textbf{Innospark-1T} & answer C \\
  & Let me analyze this question step by step. \newline The question asks: A third-year high school teacher is designing a written test to evaluate the rotational kinematics unit. If one of the evaluation indicators is "They use the mathematical relationships that describe uniform circular motion to solve various problems," which of the following items is most appropriate to evidence their achievement? \newline Let me consider each option: \newline - Option A: Determine the average angular speed of a carousel that makes 5 complete rotations in one minute \newline - Option B: Determine how many rotations and how many degrees correspond to a rotation of \textbackslash{}( \textbackslash{}frac\{5\textbackslash{}pi\}\{2\} \textbackslash{}) radians \newline - Option C: Calculate the tangential speed of a comet when its distance to the Sun is \textbackslash{}( R \textbackslash{}), if its tangential speed was 100 m/s when the distance was \textbackslash{}( \textbackslash{}frac\{R\}\{2\} \textbackslash{}) \newline \textit{[\ldots\ truncated \ldots]} \newline - Best evidences the student's ability to "use mathematical relationships to solve various problems" \newline Based on this analysis, I can determine the correct answer. \\
  \specialrule{1.2pt}{8pt}{8pt}
  \multicolumn{2}{@{}l}{\textbf{chile\_971} \quad $\Delta = -100\%$} \\
  \textbf{Question} & In a Mathematics class, the classroom team of a 3rd year Gastronomy program is working with the students on an activity of proportionality. During the monitoring, the general teacher noticed that a group of students is facing challenges in understanding proportional reasoning, specifically in the equivalence between certain values. Based on the above, which of the following strategies is pertinent to address the challenges of the students? \newline A. Pose proportion problems without numbers, using recipe ingredients as concepts \newline B. Use concrete materials visualizing the proportions of the real ingredients of a recipe by standardizing them \newline C. Integrate the proportion formula routinely, emphasizing the lost value to standardize recipes \newline D. Use value tables, observing different relationships between the number of servings and ingredients in a recipe \\
  \textbf{Gold} & D \\
  \midrule
  \textbf{Kimi-K2.6} & answer D \\
  & The user wants me to answer a multiple-choice question about teaching strategies for proportionality in a Gastronomy program. Let me analyze the scenario and options. \newline Scenario: 3rd year Gastronomy students struggling with proportional reasoning, specifically equivalence between certain values. \newline I need to find the most pertinent strategy to address this challenge. \newline Option A: Pose proportion problems without numbers, using recipe ingredients as concepts. This seems too abstract if they're already struggling with equivalence. Without numbers, it might not help with the mathematical equivalence. \newline Option B: Use concrete materials visualizing the proportions of the real ingredients of a recipe by standardizing them. This is hands-on and concrete, which is good for understanding proportions. Standardizing them helps see equivalence. In gastronomy, using real ingredients or representations could help. But "standardizing them" might mean converting to standard units, which relates to equivalence. \newline Option C: Integrate the proportion formula routinely, emphasizing the lost value to standardize recipes. "Lost value" doesn't make much sense pedagogically. Routine formula integration without conceptual understanding is not good for students struggling with reasoning. This seems like the worst option. \newline \textit{[\ldots\ truncated \ldots]} \newline Actually, looking at the options again, D seems to be the most comprehensive and appropriate pedagogical strategy. Value tables allow students to see multiple equivalent ratios and the relationship between variables. \newline I'll go with D. \\
  \textbf{Innospark-1T} & answer B \\
  & Let me analyze this question step by step. \newline The question asks: In a Mathematics class, the classroom team of a 3rd year Gastronomy program is working with the students on an activity of proportionality. During the monitoring, the general teacher noticed that a group of students is facing challenges in understanding proportional reasoning, specifically in the equivalence between certain values. Based on the abo \newline Let me consider each option: \newline - Option A: Pose proportion problems without numbers, using recipe ingredients as concepts \newline - Option B: Use concrete materials visualizing the proportions of the real ingredients of a recipe by standardizing them \newline - Option C: Integrate the proportion formula routinely, emphasizing the lost value to standardize recipes \newline \textit{[\ldots\ truncated \ldots]} \newline Final Answer: B \newline Based on this analysis, the answer is B. \\
  \specialrule{1.2pt}{8pt}{8pt}
  \multicolumn{2}{@{}l}{\textbf{chile\_655} \quad $\Delta = -100\%$} \\
  \textbf{Question} & A teacher is working with high school seniors on J. J. Rousseau's political theory, with a special emphasis on the idea of the general will, which was used by the French philosopher in his book The Social Contract (1762). Which of the following options represents a suitable example to illustrate the idea of the general will to these students? \newline A. Simple majority vote to decide the class leadership \newline B. Participation of all the students of the school and the teaching staff to reach an agreement \newline C. Request for academic reinforcement for the class, even though there are students who do not need it \newline D. Negotiation between students with opposing views to reach agreements within the class group \\
  \textbf{Gold} & C \\
  \midrule
  \textbf{Kimi-K2.6} & answer C \\
  & The user wants me to answer a multiple-choice question about Rousseau's concept of the "general will" (volont gnrale). I need to choose the best option and answer with ONLY the letter. \newline Let me recall Rousseau's concept of the general will: \newline - The general will is not simply the will of all (volont de tous), which is the sum of individual private interests. \newline - The general will aims at the common good, what is best for the community as a whole. \newline - It is not necessarily what the majority votes for at any given moment, because majority vote can reflect private interests. \newline - It requires active participation of citizens in deliberation. \newline \textit{[\ldots\ truncated \ldots]} \newline Also, consider the source context: "A teacher is working with high school seniors..." This implies the example should be pedagogically appropriate and conceptually clear. Option C is a clear, concrete example. Option B is more abstract and process-oriented. \newline Final decision: C. \\
  \textbf{Innospark-1T} & answer D \\
  & Let me analyze this question step by step. \newline The question asks: A teacher is working with high school seniors on J. J. Rousseau's political theory, with a special emphasis on the idea of the general will, which was used by the French philosopher in his book The Social Contract (1762). Which of the following options represents a suitable example to illustrate the idea of the general will to these students? \newline Let me consider each option: \newline - Option A: Simple majority vote to decide the class leadership \newline - Option B: Participation of all the students of the school and the teaching staff to reach an agreement \newline - Option C: Request for academic reinforcement for the class, even though there are students who do not need it \newline \textit{[\ldots\ truncated \ldots]} \newline While not perfect (as true general will transcends mere compromise), Option D captures the essential process Rousseau described: citizens engaging with each other to discover what truly serves their collective interest, moving beyond their individual preferences. \newline Based on this analysis, the answer is D. \\
  \specialrule{1.2pt}{8pt}{8pt}
  \multicolumn{2}{@{}l}{\textbf{chile\_416} \quad $\Delta = -100\%$} \\
  \textbf{Question} & In the unit "Interpreting urban music", a teacher asks the bassist of the course to transcribe for his instrument a score that was originally written for guitar. What prior knowledge must the student have to perform the requested transcription? \newline A. Reading bass clef notation \newline B. Reading treble clef \newline C. Reading American clef \newline D. Reading bass tablature \\
  \textbf{Gold} & B \\
  \midrule
  \textbf{Kimi-K2.6} & answer B \\
  & The question is about a bassist transcribing a score originally written for guitar. We need to determine what prior knowledge the student (bassist) must have. \newline First, let's think about standard notation for these instruments. \newline - Guitar music is typically written in treble clef (G clef), often transposed up an octave. \newline - Bass guitar (and double bass) music is typically written in bass clef (F clef). \newline The bassist needs to transcribe a guitar score for his instrument (bass). To do this, he needs to: \newline 1. Read the original guitar score (which is in treble clef). \newline \textit{[\ldots\ truncated \ldots]} \newline So B is the best answer. \newline Final answer: B. \\
  \textbf{Innospark-1T} & answer A \\
  & Let me analyze this question step by step. \newline The question asks: In the unit "Interpreting urban music", a teacher asks the bassist of the course to transcribe for his instrument a score that was originally written for guitar. What prior knowledge must the student have to perform the requested transcription? \newline Let me consider each option: \newline - Option A: Reading bass clef notation \newline - Option B: Reading treble clef \newline - Option C: Reading American clef \newline \textit{[\ldots\ truncated \ldots]} \newline The student must be able to read the source score (guitar, typically in treble clef or tablature) and write for bass (bass clef). Among the options, reading bass clef notation is essential for notating the transcribed part for bass. Therefore, the correct answer is A. \newline Based on this analysis, I can determine the correct answer. \\
  \specialrule{1.2pt}{8pt}{8pt}
  \multicolumn{2}{@{}l}{\textbf{chile\_615} \quad $\Delta = -100\%$} \\
  \textbf{Question} & In a 3rd year class, a secondary teacher is working on the following objective: "Recognize complex numbers as an extension of the numerical field of real numbers." In one of the exercises, a student presents the following procedure: 1 = -1 - -1 - -1 - -1 - i = i = -1 Most students do not find the logical problem in the procedure. Considering this evidence, what content did the students misunderstand, hindering the achievement of the stated objective? \newline A. The interpretation of the meaning of -1 \newline B. The application of the properties of roots and powers \newline C. The distinction between complex numbers that belong to the reals and those that do not \newline D. The understanding of the limitations of the properties of roots in complex numbers \\
  \textbf{Gold} & D \\
  \midrule
  \textbf{Kimi-K2.6} & answer D \\
  & The user wants me to answer a multiple-choice question about a mathematics education scenario. Let me analyze the problem carefully. \newline The scenario: \newline - Objective: "Recognize complex numbers as an extension of the numerical field of real numbers." \newline - Student's procedure: 1 = -1 - -1 - -1 - -1 - i = i = -1 \newline - Most students don't find the logical problem. \newline - Question: What content did the students misunderstand, hindering the achievement of the stated objective? \newline \textit{[\ldots\ truncated \ldots]} \newline Given the options, D is clearly the intended answer. The question is about the extension from real to complex numbers, and the limitation of root properties (specifically that (ab) = ab requires a,b 0 in the reals, and has limitations in complex numbers). \newline I'll go with D. \\
  \textbf{Innospark-1T} & answer A \\
  & Let me analyze this question step by step. \newline The question asks: In a 3rd year class, a secondary teacher is working on the following objective: "Recognize complex numbers as an extension of the numerical field of real numbers." In one of the exercises, a student presents the following procedure: 1 = -1 - -1 - -1 - -1 - -1 - i = i = -1 Most students do not find the logical p \newline Let me consider each option: \newline - Option A: The interpretation of the meaning of -1 \newline - Option B: The application of the properties of roots and powers \newline - Option C: The distinction between complex numbers that belong to the reals and those that do not \newline \textit{[\ldots\ truncated \ldots]} \newline The students misunderstood the interpretation of the meaning of \$\textbackslash{}sqrt\{-1\}\$---that it represents the imaginary unit \$i\$ in the complex number system, not an operation within the real numbers. \newline Based on this analysis, the answer is A. \\
  \specialrule{1.2pt}{8pt}{8pt}
  \multicolumn{2}{@{}l}{\textbf{chile\_621} \quad $\Delta = -100\%$} \\
  \textbf{Question} & A teacher of 8th grade, at the end of a unit, will evaluate her students' ability to select information sources, considering their reliability. To do this, she is preparing a guideline and has formulated the following evaluation indicators: 1. They identify the author of the source. 2. They briefly describe the origin or context of the source. Which of the following indicators, relevant to the content and level of the students, allows completing this evaluation guideline? \newline A. They identify the type of source \newline B. They argue about the meaning of the source \newline C. They recognize the intentionality of the source \newline D. They evaluate the validity of the data from the source \\
  \textbf{Gold} & C \\
  \midrule
  \textbf{Kimi-K2.6} & answer C \\
  & The user wants me to answer a multiple-choice question with ONLY the letter of the correct answer. No explanation. \newline Let me analyze the question: \newline Context: 8th grade teacher evaluating students' ability to select information sources considering their reliability. \newline Existing indicators: \newline 1. They identify the author of the source. \newline 2. They briefly describe the origin or context of the source. \newline \textit{[\ldots\ truncated \ldots]} \newline Also, for 8th graders, recognizing bias/intentionality is a key media literacy skill taught at that level. \newline I will go with C. \\
  \textbf{Innospark-1T} & answer A \\
  & Let me analyze this question step by step. \newline The question asks: A teacher of 8th grade, at the end of a unit, will evaluate her students' ability to select information sources, considering their reliability. To do this, she is preparing a guideline and has formulated the following evaluation indicators: 1. They identify the author of the source. 2. They briefly describe the origin or context of the source. Which of the following indic \newline Let me consider each option: \newline - Option A: They identify the type of source \newline - Option B: They argue about the meaning of the source \newline - Option C: They recognize the intentionality of the source \newline \textit{[\ldots\ truncated \ldots]} \newline Final Answer: A \newline Based on this analysis, the answer is A. \\
  \specialrule{1.2pt}{8pt}{8pt}
  \bottomrule
\end{longtable}